\section{A smooth reference join and a direct bulk response}\label{sec:response}

\subsection{A contour family that meets the reflection plane smoothly}

Let $v$ be the incoming unit ket tangent at $b$. Its reflected reverse
has outgoing tangent $-Rv$. A smooth join requires
\begin{equation}\label{eq:smooth-join}
 v=-Rv,
\end{equation}
so $v$ is parallel to the time direction $e_0$. The internal scalar
directions then agree as well, since $N(-Rv)=Mv$. For a kernel
involving many contours, every ket should have the same reference
tangent. The old semicircles arrive along $-e_3$ and return along
$+e_3$ under time reflection, producing a backtracking join instead.

Contour freedom provides a simple replacement. In the $(x^0,x^3)$-plane
set
\begin{equation}\label{eq:bulge}
 X_{T,\eta}(t)=(Tt,\eta T f(t)),\qquad f(t)=t^2(1-t),
 \quad0\leq t\leq1,\quad T>0,\quad\eta\geq0.
\end{equation}
The ket traverses the curve from $t=1$ to zero. The endpoints are on
the defect, its interior is in positive time and, for $\eta>0$,
strictly above the defect. Since $f'(0)=0$, its incoming unit tangent
is $-e_0$. The reflected return is regular with a continuous scalar
direction even when the two branches have different $T$ or $\eta$.
Curvature need not agree.

\begin{figure}[htbp]
\centering
\begin{tikzpicture}[x=1.15cm,y=1.2cm,>=Stealth]
 \fill[blue!4] (0,0) rectangle (3.7,1.75);
 \draw[->,gray] (-3.8,0)--(3.9,0) node[right] {$x^0$};
 \draw[dashed,gray] (0,-0.55)--(0,1.9) node[above] {$x^0=0$};
 \node[below] at (2,-0.45) {defect $x^3=0$};
 \node[right] at (0.08,1.7) {$x^3>0$};
 \draw[thick,MidnightBlue,domain=0:1,samples=100]
 plot ({3*\x},{8*\x*\x*(1-\x)});
 \draw[thick,BurntOrange,domain=0:1,samples=100]
 plot ({-3*\x},{8*\x*\x*(1-\x)});
 \draw[->,thick,MidnightBlue] (1.5,1)--(1.32,0.87);
 \draw[->,thick,BurntOrange] (-1.32,0.87)--(-1.5,1);
 \fill (0,0) circle (1.6pt) node[below] {$b$};
 \fill (3,0) circle (1.6pt) node[below] {$a$};
 \fill (-3,0) circle (1.6pt) node[below] {$\vartheta a$};
 \node[MidnightBlue] at (2.5,1.55) {$U_M(C)$};
 \node[BurntOrange] at (-2.5,1.55) {$U_N(\vartheta C^{-1})$};
\end{tikzpicture}
\caption{A polynomial bulge and its ordinary reflected adjoint. The
ket runs from $a$ to $b$, and the bra from $b$ to $\vartheta a$.
Both branches lie above the defect and meet with tangent $-e_0$.
The scalar maps differ away from the join. Vertical scale is enlarged.}
\label{fig:join}
\end{figure}

This construction removes a fixed-angle geometric and scalar-direction
cusp in the direct bulk term. It does not prove absence of defect
contact or junction counterterms. No driver or capacity law for
\eqref{eq:bulge} is computed here; it is a smooth contour control for
the field-theory problem, not an identification with $\omega$.

\subsection{The specified Gaussian perturbative sector}

Use the free Feynman-gauge propagators
\begin{equation}\label{eq:bulk-propagators}
 \vev{A_\mu^a(x)A_\nu^b(y)}_0
 =\frac{g^2\delta^{ab}\delta_{\mu\nu}}{4\pi^2|x-y|^2},\qquad
 \vev{X_A^a(x)X_B^b(y)}_0
 =\frac{g^2\delta^{ab}\delta_{AB}}{4\pi^2|x-y|^2}.
\end{equation}
This defines the direct exchange of a bulk field between the two line
branches. Endpoint self-energy, endpoint-to-line interactions, defect
vertices and counterterms are excluded from this sector. Only their
consistent sum could be a gauge-independent interacting answer.

On a single branch the gauge/scalar numerator cancels by $M^TM=I$
or $N^TN=I$. If $v_i,v_j$ denote oriented original ket tangents,
the cross-branch numerator instead is
\begin{align}\label{eq:cross-numerator}
 -(-Rv_i)\cdot v_j+[N(-Rv_i)]\cdot Mv_j
 &=v_i\cdot(R+I)v_j\notag\\
 &=2v_{i,\mathrm{sp}}\cdot v_{j,\mathrm{sp}}.
\end{align}
Here $\mathrm{sp}$ consists of the three directions unchanged by time
reflection, including the defect normal. Near a regular reference
join satisfying \eqref{eq:smooth-join}, the spatial tangents vanish
linearly in arclength. The numerator is $O(st)$ and the denominator
is comparable to $(s+t)^2$. The local double integral is finite.
This local estimate concerns \eqref{eq:bulk-propagators}, not all
possible defect graphs.

For the self-pairing of \eqref{eq:bulge}, $T$ cancels by scale
invariance and the stripped cross term is
\begin{equation}\label{eq:bulk-I}
 I(\eta)=2\eta^2\int_0^1\!\int_0^1
 \frac{f'(s)f'(t)}{(s+t)^2+\eta^2[f(s)-f(t)]^2}\dd s\dd t.
\end{equation}
If $\sum_aT_aT_a=C_F I$ in the endpoint representation, its contribution
relative to the free endpoint propagator is
$g^2C_F I(\eta)/(4\pi^2)$. There is no factor $1/2$ for an exchange
between two distinct ordered branches.

\begin{proposition}[The direct bulk bulge coefficient]\label{prop:bulge-coefficient}
For $f(t)=t^2(1-t)$,
\begin{equation}\label{eq:bulge-coefficient}
 I(\eta)=(\tfrac32-2\log2)\eta^2+O(\eta^4),\qquad
 \tfrac32-2\log2=0.1137056388801094\ldots>0.
\end{equation}
One may bound the absolute remainder by $2\eta^4$ for real $\eta$.
\end{proposition}
\begin{proof}
The quadratic coefficient is $2A$, where
\begin{equation}
 A=\int_0^1\!\int_0^1
 \frac{(2s-3s^2)(2t-3t^2)}{(s+t)^2}\dd s\dd t.
\end{equation}
Put $r=s/(s+t)$ and $u=s+t$. The Jacobian is $u$.
By symmetry, integrate $0\leq r\leq1/2$ and
$0\leq u\leq1/(1-r)$ and multiply by two. The integrand including
the Jacobian is
$r(1-r)[4-6u+9r(1-r)u^2]u$.
Its $u$ integral, including that factor two, simplifies to
$r^2/[2(1-r)^2]$. Hence
\begin{equation}
 A=\frac12\int_0^{1/2}\frac{r^2}{(1-r)^2}\dd r
 =\frac34-\log2.
\end{equation}
For the remainder, $|f'|\leq1$, $|f'(s)|\leq2s$ and
$|f(s)-f(t)|\leq|s-t|\leq s+t$. Subtracting the denominator at
$\eta=0$ bounds the absolute integrand difference by
\begin{equation}
 2\eta^4\frac{|f'(s)f'(t)|[f(s)-f(t)]^2}{(s+t)^4}
 \leq8\eta^4\frac{st}{(s+t)^2}\leq2\eta^4.
\end{equation}
The integration square has area one. This justifies the expansion
uniformly even at the apparent singular corner.
\end{proof}

Numerical controls give $I(0)=0$,
$I(0.3)\simeq0.0102220294386856$ and
$I(0.7)\simeq0.0553790171218514$. These check the integral and its
small-bulge coefficient; they are not estimates of the complete
interacting response.

There is also a positive Gaussian interpretation of this cross term.
The reflected four-dimensional scalar covariance has representation
\begin{equation}
 D_0(\vartheta x-y)=\int_{\R^3}\frac{\mathrm d^3k}{(2\pi)^3}
 \frac{e^{-|k|(x^0+y^0)}e^{\ii k\cdot(\mathbf y-\mathbf x)}}{2|k|}.
\end{equation}
Smear it against the three spatial tangent currents of each ket and
sum the components as in \eqref{eq:cross-numerator}. This is a Gram
matrix. The vanishing of those currents at $b$ permits the finite
join limit just estimated. This proves positivity of the displayed
Gaussian current sector, without extending it to the omitted diagrams.

\subsection{Why a positive norm need not be protected}

For an actual reflected pairing a smooth deformation has the structure
\begin{equation}
 \delta K_{ij}=\vev{\Theta_0(\delta F_i)F_j}
                 +\vev{\Theta_0F_i\,\delta F_j},
\end{equation}
with the insertions of Section~\ref{sec:variation} and any renormalized
contact terms. Even if a deformation of a ket is $Q$-exact,
$\delta\psi=Q\lambda$, and $Q\psi=0$, its physical norm obeys
\begin{equation}\label{eq:Q-norm}
 \delta\norm\psi^2=2\Ree\ip\psi{Q\lambda}
                  =2\Ree\ip{Q^\dagger\psi}\lambda.
\end{equation}
The first condition does not imply $Q^\dagger\psi=0$. For example,
on a Hilbert space with orthonormal $e_1,e_2$, let
$Qe_2=e_1$, $Qe_1=0$, $\psi=e_1$, $\lambda=e_2$.
Then $Q^2=0$, $Q\psi=0$, but \eqref{eq:Q-norm} equals two.
Adding a vector $e_3$ killed by both $Q$ and $Q^\dagger$ supplies a
zero-first-variation control. Thus a common-charge Ward identity
and a physical positive norm are different requirements.
